\documentclass[11pt,a4paper]{article}
\usepackage[utf8]{inputenc}
\usepackage[T1]{fontenc}
\usepackage[brazilian]{babel}
\usepackage{amsmath,amssymb,amsthm}
\usepackage{geometry}
\geometry{top=2.2cm,bottom=2.2cm,left=2cm,right=2cm}
\usepackage{booktabs}
\usepackage{array}
\usepackage{xcolor}
\usepackage{fancyhdr}
\usepackage{tcolorbox}
\tcbuselibrary{skins,breakable}
\usepackage{tikz}
\usetikzlibrary{trees,arrows.meta,positioning}
\usepackage{listings}
\usepackage{hyperref}
\hypersetup{
    colorlinks=true,
    linkcolor=blue!75!black,
    urlcolor=blue!75!black,
    citecolor=blue!75!black
}

\setlength{\headheight}{14.5pt}
\pagestyle{fancy}
\fancyhf{}
\lhead{\textbf{UFU -- FEELT} | PAA: Guia Didático de Estudos}
\rhead{Prof. Felipe A. Louza | \textbf{Lista 9}}
\lfoot{\small Arthur Pires Lopes -- 12221EEL001}
\rfoot{\thepage}

\setlength{\parindent}{0pt}
\setlength{\parskip}{5pt}

% Estilos de caixas didáticas (sem emojis unicode para compatibilidade total com pdflatex)
\newtcolorbox{dica}[1][]{
    enhanced,
    colback=blue!5!white,
    colframe=blue!75!black,
    fonttitle=\bfseries\small,
    title={[Dica de Estudo / Ponto de Atenção] #1},
    boxrule=0.8pt,
    arc=3pt,
    breakable
}

\newtcolorbox{teoria}[1][]{
    enhanced,
    colback=teal!5!white,
    colframe=teal!80!black,
    fonttitle=\bfseries\small,
    title={[Fundamentação Teórica] #1},
    boxrule=0.8pt,
    arc=3pt,
    breakable
}

\newtcolorbox{alerta}[1][]{
    enhanced,
    colback=red!5!white,
    colframe=red!75!black,
    fonttitle=\bfseries\small,
    title={[Atenção: Pegadinha Clássica de Prova] #1},
    boxrule=0.8pt,
    arc=3pt,
    breakable
}

\lstset{
  basicstyle=\ttfamily\small,
  frame=single,
  rulecolor=\color{black!30},
  backgroundcolor=\color{black!3},
  numbers=left,
  numberstyle=\tiny\color{gray},
  stepnumber=1,
  numbersep=8pt,
  showstringspaces=false,
  breaklines=true,
  literate={á}{{\'a}}1 {ã}{{\~a}}1 {ç}{{\c c}}1 {é}{{\'e}}1 {í}{{\'i}}1 {ó}{{\'o}}1 {ú}{{\'u}}1 {Á}{{\'A}}1 {É}{{\'E}}1 {Í}{{\'I}}1 {Ó}{{\'O}}1 {Ú}{{\'U}}1 {≤}{{$\le$}}1 {≥}{{$\ge$}}1 {←}{{$\leftarrow$}}1 {≠}{{$\ne$}}1 {∞}{{$\infty$}}1
}

\begin{document}

% =========================================================================
% CABEÇALHO DO GUIA DE ESTUDOS
% =========================================================================
\begin{center}
    {\Large \textbf{Universidade Federal de Uberlândia -- UFU}}\\[2pt]
    {\large Faculdade de Engenharia Elétrica (FEELT)}\\[2pt]
    {\large Projeto e Análise de Algoritmos (PAA) -- Prof. Felipe A. Louza}\\[4pt]
    \rule{\textwidth}{1pt}\\[5pt]
    {\LARGE \textbf{Lista 9: Programação Dinâmica}}\\[3pt]
    {\large \textbf{Guia Completo de Estudos e Resoluções Comentadas}}\\[3pt]
    \textbf{Foco: Prova 2 e Concursos de Alto Nível}\\[2pt]
    \rule{\textwidth}{0.5pt}
\end{center}

\vspace{0.1cm}
\noindent \textbf{Aluno:} Arthur Pires Lopes \quad \textbf{Matrícula:} 12221EEL001

\vspace{0.3cm}

\begin{teoria}[Os Dois Pilares Fundamentais da Programação Dinâmica (PD)]
Para que um problema de otimização possa ser resolvido eficientemente por Programação Dinâmica, ele deve satisfazer duas propriedades essenciais:
\begin{enumerate}
    \item \textbf{Subestrutura Ótima:} Uma solução ótima para o problema global é composta por soluções ótimas para seus subproblemas. Em termos matemáticos, se tomamos uma decisão local ótima (ex: primeiro corte, inclusão ou não de um item), o problema restante deve ser solucionado de forma ótima independentemente.
    \item \textbf{Sobreposição de Subproblemas (Overlapping Subproblems):} Um algoritmo recursivo ingênuo visitaria os mesmos subproblemas uma quantidade exponencial de vezes (árvore de recursão redundante). Em contrapartida, o espaço total de subproblemas distintos é pequeno (polinomial no tamanho da entrada, como $O(n)$, $O(n^2)$ ou $O(n \cdot W)$). A PD salva as soluções calculadas em uma tabela (vetor ou matriz), evitando qualquer recomputação (\emph{troca tempo por espaço}).
\end{enumerate}
A dependência entre os subproblemas forma um \textbf{Grafo Acíclico Dirigido (DAG)}. A estratégia \emph{Bottom-Up} resolve os subproblemas na \textbf{ordenação topológica reversa} desse DAG (do caso base para o caso geral).
\end{teoria}

\vspace{0.4cm}

% =========================================================================
% QUESTÃO 1
% =========================================================================
\section*{Questão 1: Corte de Barras (Cut-Rod Bottom-Up)}

\textbf{Enunciado Oficial:} Mostre a execução passo-a-passo do algoritmo \textsc{Cut-Rod-Bottom-Up}$(P, n)$ para a seguinte entrada:
\begin{center}
\begin{tabular}{r|ccccccc}
\toprule
\textbf{tamanho} $i$ & 0 & 1 & 2 & 3 & 4 & 5 & 6 \\
\midrule
\textbf{preço} $p_i$ & 0 & 1 & 3 & 11 & 16 & 19 & 10 \\
\bottomrule
\end{tabular}
\end{center}

\begin{teoria}[Equação de Bellman para o Corte de Barras]
Seja $A[i]$ a receita máxima obtida para uma barra de comprimento $i$. Ao realizarmos um corte de tamanho $j \in \{1, \dots, i\}$, vendemos o primeiro pedaço por $p_j$ e ficamos com uma barra restante de tamanho $i - j$, cuja solução ótima já foi calculada e vale $A[i - j]$.
\[
A[0] = 0, \qquad A[i] = \max_{1 \le j \le i} \Big( p_j + A[i - j] \Big) \quad \text{para } i \ge 1.
\]
\end{teoria}

\begin{dica}[Mecânica Prática para a Prova]
Para preencher a posição $A[i]$, basta somar os valores das extremidades opostas:
\[
p_1 + A[i-1], \quad p_2 + A[i-2], \quad p_3 + A[i-3], \quad \dots, \quad p_i + A[0].
\]
O maior valor entre essas somas é o resultado que entra em $A[i]$.
\end{dica}

\subsection*{Execução Passo a Passo Detalhada}

\begin{itemize}
    \item \textbf{Base $i = 0$:} Uma barra de comprimento 0 não rende nada: $A[0] = 0$.
    \item \textbf{Comprimento $i = 1$:} Única opção é vender o pedaço inteiro de 1:
    \[
    j = 1: p_1 + A[0] = 1 + 0 = 1 \implies A[1] = 1.
    \]
    \item \textbf{Comprimento $i = 2$:}
    \begin{align*}
    j = 1 &: p_1 + A[1] = 1 + 1 = 2 \quad (\text{dois pedaços de 1}) \\
    j = 2 &: p_2 + A[0] = 3 + 0 = 3 \quad (\text{um pedaço inteiro de 2})
    \end{align*}
    $\implies A[2] = \max(2, 3) = 3$ (com corte em $j=2$).
    \item \textbf{Comprimento $i = 3$:}
    \begin{align*}
    j = 1 &: p_1 + A[2] = 1 + 3 = 4 \\
    j = 2 &: p_2 + A[1] = 3 + 1 = 4 \\
    j = 3 &: p_3 + A[0] = 11 + 0 = 11 \quad (\text{muito lucrativo!})
    \end{align*}
    $\implies A[3] = \max(4, 4, 11) = 11$ (com corte em $j=3$).
    \item \textbf{Comprimento $i = 4$:}
    \begin{align*}
    j = 1 &: p_1 + A[3] = 1 + 11 = 12 \\
    j = 2 &: p_2 + A[2] = 3 + 3 = 6 \\
    j = 3 &: p_3 + A[1] = 11 + 1 = 12 \\
    j = 4 &: p_4 + A[0] = 16 + 0 = 16
    \end{align*}
    $\implies A[4] = \max(12, 6, 12, 16) = 16$ (com corte em $j=4$).
    \item \textbf{Comprimento $i = 5$:}
    \begin{align*}
    j = 1 &: p_1 + A[4] = 1 + 16 = 17 \\
    j = 2 &: p_2 + A[3] = 3 + 11 = 14 \\
    j = 3 &: p_3 + A[2] = 11 + 3 = 14 \\
    j = 4 &: p_4 + A[1] = 16 + 1 = 17 \\
    j = 5 &: p_5 + A[0] = 19 + 0 = 19
    \end{align*}
    $\implies A[5] = \max(17, 14, 14, 17, 19) = 19$ (com corte em $j=5$).
    \item \textbf{Comprimento $i = 6$:}
    \begin{align*}
    j = 1 &: p_1 + A[5] = 1 + 19 = 20 \\
    j = 2 &: p_2 + A[4] = 3 + 16 = 19 \\
    j = 3 &: p_3 + A[3] = 11 + 11 = \mathbf{22} \quad \bigstar \textbf{ MÁXIMO} \\
    j = 4 &: p_4 + A[2] = 16 + 3 = 19 \\
    j = 5 &: p_5 + A[1] = 19 + 1 = 20 \\
    j = 6 &: p_6 + A[0] = 10 + 0 = 10 \quad (\text{não cortar dá prejuízo: apenas 10})
    \end{align*}
    $\implies A[6] = \max(20, 19, \mathbf{22}, 19, 20, 10) = \mathbf{22}$.
\end{itemize}

\begin{alerta}[Atenção à Desproporção de Preços na Prova]
Note que $p_6 = 10$, enquanto $p_3 = 11$. Uma barra inteira de tamanho 6 vale menos do que uma única barra de tamanho 3! Por isso, cortar a barra de tamanho 6 em dois pedaços de tamanho 3 ($3 + 3$) gera $11 + 11 = 22$, mais que o dobro do preço da barra intacta ($10$). A Programação Dinâmica captura essa não-linearidade perfeitamente.
\end{alerta}

\textbf{Tabela Final Resumo:}
\begin{center}
\begin{tabular}{r|ccccccc}
\toprule
$i$ & 0 & 1 & 2 & 3 & 4 & 5 & 6 \\
\midrule
$p_i$ & 0 & 1 & 3 & 11 & 16 & 19 & 10 \\
$A[i]$ & \textbf{0} & \textbf{1} & \textbf{3} & \textbf{11} & \textbf{16} & \textbf{19} & \textbf{22} \\
\bottomrule
\end{tabular}
\end{center}

\vspace{0.4cm}

% =========================================================================
% QUESTÃO 2
% =========================================================================
\section*{Questão 2: Reconstrução dos Cortes Ótimos}

\textbf{Enunciado Oficial:} Escreva um algoritmo para descobrir os tamanhos dos pedaços nos quais a barra foi cortada na solução ótima. \emph{Dica: Modifique o algoritmo \textsc{Cut-Rod-Bottom-Up}$(P, n)$ para registrar os cortes na solução.}

\begin{teoria}[Padrão de Reconstrução em Programação Dinâmica]
A tabela DP armazena apenas o \textbf{valor escalar} da solução ótima ($A[n]$). Para resgatar a \textbf{estrutura da solução} (quais pedaços foram cortados), mantemos uma tabela de decisões $S[1 \dots n]$. 
Quando $p_j + A[i-j]$ supera o máximo provisório, salvamos $S[i] = j$. Posteriormente, retrocedemos iterativamente: a barra tem tamanho $n$, extraímos o primeiro pedaço $S[n]$, e o problema se reduz para uma barra de tamanho $n - S[n]$.
\end{teoria}

\subsection*{Algoritmos Completos e Comentados}

\begin{lstlisting}[caption={Algoritmo com Registro de Cortes Auxiliares}]
EXTENDED-CUT-ROD-BOTTOM-UP(P, n)
1  let A[0 .. n] and S[1 .. n] be new arrays
2  A[0] = 0
3  for i = 1 to n
4      m = -infinity
5      for j = 1 to i
6          if m < P[j] + A[i - j]
7              m = P[j] + A[i - j]
8              S[i] = j             // Registra que o melhor corte inicial foi j
9      A[i] = m
10 return A and S
\end{lstlisting}

\begin{lstlisting}[caption={Algoritmo de Recuperação e Retorno dos Pedaços}]
RETRIEVE-CUT-ROD-SOLUTION(P, n)
1  (A, S) = EXTENDED-CUT-ROD-BOTTOM-UP(P, n)
2  let L be an empty list
3  curr = n
4  while curr > 0
5      add S[curr] to L             // Adiciona o tamanho do pedaco na lista
6      curr = curr - S[curr]        // Reduz a barra pelo tamanho cortado
7  return L
\end{lstlisting}

\subsection*{Rastreamento com o Exemplo da Questão 1}
Para a entrada dada:
\begin{center}
\begin{tabular}{r|cccccc}
\toprule
$i$ & 1 & 2 & 3 & 4 & 5 & 6 \\
\midrule
$S[i]$ & 1 & 2 & 3 & 4 & 5 & \textbf{3} \\
\bottomrule
\end{tabular}
\end{center}
\begin{enumerate}
    \item Começamos com $curr = 6$. Consultamos $S[6] = 3 \implies$ Cortamos um pedaço de $3$.
    \item Nova barra restante: $curr = 6 - 3 = 3$.
    \item Consultamos $S[3] = 3 \implies$ Cortamos outro pedaço de $3$.
    \item Nova barra restante: $curr = 3 - 3 = 0$. Condição de parada atingida.
    \item Lista retornada: $\langle 3, 3 \rangle$. Soma dos tamanhos: $3 + 3 = 6$. Receita: $p_3 + p_3 = 11 + 11 = 22$.
\end{enumerate}

\begin{dica}[Complexidade para Provas]
\begin{itemize}
    \item \textbf{Tempo de Preenchimento:} $\sum_{i=1}^n i = \frac{n(n+1)}{2} = \Theta(n^2)$.
    \item \textbf{Tempo de Reconstrução:} O laço \texttt{while} executa no máximo $n$ passos (pois $S[curr] \ge 1$ sempre), logo leva tempo $O(n)$.
    \item \textbf{Tempo Total:} $\Theta(n^2)$.
    \item \textbf{Espaço Auxiliar:} $\Theta(n)$ para os vetores $A$ e $S$.
\end{itemize}
\end{dica}

\vspace{0.4cm}

% =========================================================================
% QUESTÃO 3
% =========================================================================
\section*{Questão 3: Maior Subsequência Crescente (LIS)}

\textbf{Enunciado Oficial:} Mostre a execução passo-a-passo do algoritmo \textsc{LIS-PD}$(S, n)$ para a seguinte entrada:
\begin{center}
\begin{tabular}{r|cccccccc}
\toprule
$i$ & 1 & 2 & 3 & 4 & 5 & 6 & 7 & 8 \\
\midrule
$s_i$ & 5 & 2 & 8 & 6 & 3 & 6 & 9 & 7 \\
\bottomrule
\end{tabular}
\end{center}

\begin{teoria}[Definição Formal do Subproblema da LIS]
Para uma sequência $S = \langle s_1, s_2, \dots, s_n \rangle$, definimos $L[i]$ como o comprimento da maior subsequência estritamente crescente que \textbf{obrigatoriamente termina no elemento $s_i$}.
A relação de recorrência é:
\[
L[1] = 1, \qquad L[i] = 1 + \max \left( \{ L[j] \mid 1 \le j < i \text{ e } s_j < s_i \} \cup \{0\} \right).
\]
O comprimento da LIS global é o valor máximo dentre todos os términos possíveis:
\[
\text{LIS}(S) = \max_{1 \le j \le n} L[j].
\]
\end{teoria}

\begin{alerta}[Pegadinha Mortal: Crescimento Estrito vs. Não-Estrito]
O enunciado exige elementos \textbf{estritamente crescentes} ($s_j < s_i$). Se $s_j = s_i$, o elemento anterior $s_j$ \textbf{não pode} estender a subsequência!
Veja o caso $i = 6$ ($s_6 = 6$): os elementos anteriores são $s_1=5, s_2=2, s_3=8, s_4=6, s_5=3$. O elemento $s_4 = 6$ \textbf{não é menor} que $s_6 = 6$ ($6 \not< 6$). Portanto, $L[4]$ não entra na disputa de candidatos para $L[6]$!
\end{alerta}

\subsection*{Cálculos Célula a Célula}

\begin{itemize}
    \item \textbf{$i = 1$ ($s_1 = 5$):} Primeiro elemento, não há anteriores. $L[1] = 1$. LIS: $\langle 5 \rangle$.
    \item \textbf{$i = 2$ ($s_2 = 2$):}
    $j=1: s_1=5 \not< 2$. Candidatos vazios $\implies L[2] = 1 + 0 = 1$. LIS: $\langle 2 \rangle$.
    \item \textbf{$i = 3$ ($s_3 = 8$):}
    \begin{itemize}
        \item $j=1: s_1=5 < 8 \implies L[1] = 1$
        \item $j=2: s_2=2 < 8 \implies L[2] = 1$
    \end{itemize}
    Maior anterior $\ell = 1 \implies L[3] = 1 + 1 = 2$. LIS: $\langle 5, 8 \rangle$ ou $\langle 2, 8 \rangle$.
    \item \textbf{$i = 4$ ($s_4 = 6$):}
    \begin{itemize}
        \item $j=1: s_1=5 < 6 \implies L[1] = 1$
        \item $j=2: s_2=2 < 6 \implies L[2] = 1$
        \item $j=3: s_3=8 \not< 6$
    \end{itemize}
    Maior anterior $\ell = 1 \implies L[4] = 1 + 1 = 2$. LIS: $\langle 5, 6 \rangle$ ou $\langle 2, 6 \rangle$.
    \item \textbf{$i = 5$ ($s_5 = 3$):}
    \begin{itemize}
        \item $j=1: s_1=5 \not< 3$
        \item $j=2: s_2=2 < 3 \implies L[2] = 1$
        \item $j=3, 4: s_3=8 \not< 3, s_4=6 \not< 3$
    \end{itemize}
    Maior anterior $\ell = 1 \implies L[5] = 1 + 1 = 2$. LIS: $\langle 2, 3 \rangle$.
    \item \textbf{$i = 6$ ($s_6 = 6$):}
    \begin{itemize}
        \item $j=1: s_1=5 < 6 \implies L[1] = 1$
        \item $j=2: s_2=2 < 6 \implies L[2] = 1$
        \item $j=3: s_3=8 \not< 6$
        \item $j=4: s_4=6 \not< 6$ (iguais, não estritamente menor!)
        \item $j=5: s_5=3 < 6 \implies L[5] = 2$
    \end{itemize}
    Maior anterior $\ell = \max(1, 1, 2) = 2 \implies L[6] = 1 + 2 = \mathbf{3}$. LIS: $\langle 2, 3, 6 \rangle$.
    \item \textbf{$i = 7$ ($s_7 = 9$):}
    Candidatos menores que 9: $s_1=5(L=1), s_2=2(L=1), s_3=8(L=2), s_4=6(L=2), s_5=3(L=2), s_6=6(L=3)$.
    Maior anterior $\ell = \max(1, 1, 2, 2, 2, 3) = 3 \implies L[7] = 1 + 3 = \mathbf{4}$. LIS: $\langle 2, 3, 6, 9 \rangle$.
    \item \textbf{$i = 8$ ($s_8 = 7$):}
    Candidatos menores que 7: $s_1=5(L=1), s_2=2(L=1), s_4=6(L=2), s_5=3(L=2), s_6=6(L=3)$.
    (Excluídos: $s_3 = 8 \not< 7$ e $s_7 = 9 \not< 7$).
    Maior anterior $\ell = \max(1, 1, 2, 2, 3) = 3 \implies L[8] = 1 + 3 = \mathbf{4}$. LIS: $\langle 2, 3, 6, 7 \rangle$.
\end{itemize}

\textbf{Tabela Completa de Resultados:}
\begin{center}
\begin{tabular}{r|cccccccc}
\toprule
$i$ & 1 & 2 & 3 & 4 & 5 & 6 & 7 & 8 \\
\midrule
$s_i$ & 5 & 2 & 8 & 6 & 3 & 6 & 9 & 7 \\
$L[i]$ & \textbf{1} & \textbf{1} & \textbf{2} & \textbf{2} & \textbf{2} & \textbf{3} & \textbf{4} & \textbf{4} \\
\bottomrule
\end{tabular}
\end{center}

\textbf{Resposta Final:} A subsequência crescente mais longa tem tamanho $\max_j L[j] = \mathbf{4}$.

\vspace{0.4cm}

% =========================================================================
% QUESTÃO 4
% =========================================================================
\section*{Questão 4: Reconstrução da Sequência LIS Completa}

\textbf{Enunciado Oficial:} Escreva um algoritmo que retorna a LIS de uma sequência (não apenas o tamanho da LIS). \emph{Dica: Modifique o algoritmo \textsc{LIS-PD}$(S, n)$ para registrar as soluções parciais do problema.}

\begin{teoria}[Vetor de Predecessores para a LIS]
Assim como no cálculo de caminhos mínimos em grafos, a LIS constrói implicitamente uma árvore de predecessores. Para recuperar a lista de elementos na ordem correta:
\begin{enumerate}
    \item Criamos um vetor $P[1 \dots n]$, com $P[i] = -1$ inicialmente.
    \item Ao calcular $L[i]$, se um elemento $j < i$ é o predecessor ótimo ($s_j < s_i$ e maximiza $L[j]$), definimos $P[i] = j$.
    \item Ao final, tomamos $k = \operatorname{argmax}_j L[j]$ e seguimos os ponteiros: $k \to P[k] \to P[P[k]] \dots$ até $-1$.
    \item A lista coletada estará invertida (do último para o primeiro). Invertendo-a, obtemos a LIS ótima.
\end{enumerate}
\end{teoria}

\subsection*{Algoritmo Completo}

\begin{lstlisting}[caption={Algoritmo LIS com Retorno da Sequência Completa}]
LIS-COMPLETA(S, n)
1  let L[1 .. n] and P[1 .. n] be new arrays
2  L[1] = 1
3  P[1] = -1
4  max_len = 1
5  best_idx = 1
6  for i = 2 to n
7      l = 0
8      p = -1
9      for j = 1 to i - 1
10         if S[j] < S[i] and L[j] > l
11             l = L[j]
12             p = j
13     L[i] = 1 + l
14     P[i] = p
15     if L[i] > max_len
16         max_len = L[i]
17         best_idx = i
18 // Reconstrucao por retrocesso
19 let R be an empty list
20 curr = best_idx
21 while curr != -1
22     add S[curr] to R
23     curr = P[curr]
24 reverse(R)
25 return R
\end{lstlisting}

\begin{dica}[Execução com a Questão 3]
Para o exemplo com $best\_idx = 7$ ($s_7 = 9$):
$P[7] = 6$ ($s_6=6$), $P[6] = 5$ ($s_5=3$), $P[5] = 2$ ($s_2=2$), $P[2] = -1$.
A lista retrocedida é $\langle 9, 6, 3, 2 \rangle$. Invertida: $\langle 2, 3, 6, 9 \rangle$, de tamanho 4.
\end{dica}

\vspace{0.4cm}

% =========================================================================
% QUESTÃO 5
% =========================================================================
\section*{Questão 5: Problema do Troco (Coin Change)}

\textbf{Enunciado Oficial:} Mostre a execução passo-a-passo do algoritmo \textsc{COINS-PD}$(C, k)$ para a seguinte entrada:
\[
C = [1, 5, 8, 11] \quad \text{e} \quad k = 13
\]

\begin{teoria}[Equação de Bellman para o Problema do Troco Clássico]
Seja $M[i]$ o número mínimo de moedas para totalizar o valor $i$. Como temos quantidade ilimitada de cada tipo de moeda:
\[
M[0] = 0, \qquad M[i] = 1 + \min_{\substack{1 \le j \le n \\ i - c_j \ge 0}} M[i - c_j].
\]
\end{teoria}

\subsection*{Preenchimento Passo a Passo do Vetor $M[0 \dots 13]$}
\begin{itemize}
    \item $i = 0$: $M[0] = 0$ (quantia zero exige zero moedas).
    \item $i = 1$: apenas moeda 1 cabe: $1 + M[0] = 1 \implies M[1] = 1$.
    \item $i = 2$: $1 + M[1] = 2 \implies M[2] = 2$.
    \item $i = 3$: $1 + M[2] = 3 \implies M[3] = 3$.
    \item $i = 4$: $1 + M[3] = 4 \implies M[4] = 4$.
    \item $i = 5$: cabem moedas 1 e 5:
    $1 + \min(M[4], M[0]) = 1 + \min(4, 0) = 1 \implies M[5] = 1$. (1 moeda de 5)
    \item $i = 6$: cabem moedas 1 e 5:
    $1 + \min(M[5], M[1]) = 1 + \min(1, 1) = 2 \implies M[6] = 2$. (moedas $5+1$)
    \item $i = 7$: cabem moedas 1 e 5:
    $1 + \min(M[6], M[2]) = 1 + \min(2, 2) = 3 \implies M[7] = 3$. (moedas $5+1+1$)
    \item $i = 8$: cabem moedas 1, 5 e 8:
    $1 + \min(M[7], M[3], M[0]) = 1 + \min(3, 3, 0) = 1 \implies M[8] = 1$. (1 moeda de 8)
    \item $i = 9$: cabem moedas 1, 5 e 8:
    $1 + \min(M[8], M[4], M[1]) = 1 + \min(1, 4, 1) = 2 \implies M[9] = 2$. (moedas $8+1$)
    \item $i = 10$: cabem moedas 1, 5 e 8:
    $1 + \min(M[9], M[5], M[2]) = 1 + \min(2, 1, 2) = 2 \implies M[10] = 2$. (moedas $5+5$)
    \item $i = 11$: cabem moedas 1, 5, 8 e 11:
    $1 + \min(M[10], M[6], M[3], M[0]) = 1 + \min(2, 2, 3, 0) = 1 \implies M[11] = 1$. (1 moeda de 11)
    \item $i = 12$: cabem moedas 1, 5, 8 e 11:
    $1 + \min(M[11], M[7], M[4], M[1]) = 1 + \min(1, 3, 4, 1) = 2 \implies M[12] = 2$. (moedas $11+1$)
    \item $i = 13$: cabem moedas 1, 5, 8 e 11:
    \begin{align*}
    M[13] &= 1 + \min\Big( M[13-1], M[13-5], M[13-8], M[13-11] \Big) \\
          &= 1 + \min\Big( M[12], M[8], M[5], M[2] \Big) \\
          &= 1 + \min\Big( 2, 1, 1, 2 \Big) = 1 + 1 = \mathbf{2}.
    \end{align*}
\end{itemize}

\begin{alerta}[Por que o Algoritmo Guloso Falharia?]
Um algoritmo guloso sempre escolhe a maior moeda possível. Para $k = 13$, o guloso pegaria a moeda de $11$, restando $13 - 11 = 2$, precisando então de duas moedas de $1$ ($11 + 1 + 1 = 13$), totalizando \textbf{3 moedas}.
A Programação Dinâmica avalia todos os subproblemas e descobre a combinação ótima: $8 + 5 = 13$, que usa apenas \textbf{2 moedas}! Essa é uma das razões clássicas pelas quais PD é ensinada após algoritmos gulosos.
\end{alerta}

\textbf{Tabela Completa de $M$:}
\begin{center}
\begin{tabular}{r|cccccccccccccc}
\toprule
$i$ & 0 & 1 & 2 & 3 & 4 & 5 & 6 & 7 & 8 & 9 & 10 & 11 & 12 & 13 \\
\midrule
$M[i]$ & \textbf{0} & \textbf{1} & \textbf{2} & \textbf{3} & \textbf{4} & \textbf{1} & \textbf{2} & \textbf{3} & \textbf{1} & \textbf{2} & \textbf{2} & \textbf{1} & \textbf{2} & \textbf{2} \\
\bottomrule
\end{tabular}
\end{center}

\vspace{0.4cm}

% =========================================================================
% QUESTÃO 6
% =========================================================================
\section*{Questão 6: Troco sem Repetição de Moedas (0/1 Coin Change)}

\textbf{Enunciado Oficial:} Escreva um algoritmo que resolve o problema das moedas no qual \textbf{não é permitido repetições de moedas no troco}.

\begin{teoria}[Mapeamento para o Problema da Mochila 0-1]
Quando cada moeda $c_i \in C$ é um recurso finito (disponível apenas 1 unidade), temos uma restrição $0/1$.
Definimos $M[i, v]$ como o número mínimo de moedas dentre $\{c_1, \dots, c_i\}$ para formar a quantia exata $v$.
\begin{itemize}
    \item \textbf{Opção 1 (Não usar a moeda $i$):} O troco $v$ deve ser formado exclusivamente pelas moedas anteriores: $M[i-1, v]$.
    \item \textbf{Opção 2 (Usar a moeda $i$ uma única vez):} Consumimos 1 moeda de valor $c_i$, e o valor restante $v - c_i$ deve vir das $i-1$ moedas anteriores: $1 + M[i-1, v - c_i]$.
\end{itemize}
Recorrência formal:
\[
M[i, v] = 
\begin{cases}
0, & \text{se } v = 0 \\
\infty, & \text{se } i = 0 \text{ e } v > 0 \\
M[i-1, v], & \text{se } c_i > v \\
\min\Big( M[i-1, v], \; 1 + M[i-1, v - c_i] \Big), & \text{se } c_i \le v
\end{cases}
\]
\end{teoria}

\subsection*{Algoritmo em Pseudocódigo}

Abaixo apresentamos a versão elegante com vetor unidimensional $O(k)$ de espaço.

\begin{lstlisting}[caption={Algoritmo do Troco sem Repetição (0/1)}]
COINS-NO-REPEAT(C, n, k)
1  let dp[0 .. k] be a new array
2  dp[0] = 0
3  for v = 1 to k
4      dp[v] = infinity
5  for i = 1 to n
6      for v = k down to C[i]       // ITERACAO DECRESCENTE E OBRIGATORIA!
7          if dp[v - C[i]] != infinity
8              dp[v] = min(dp[v], 1 + dp[v - C[i]])
9  if dp[k] == infinity
10     return -1    // Impossivel formar o valor k sem repetir moedas
11 return dp[k]
\end{lstlisting}

\begin{alerta}[Por que o laço decrescente na linha 6 é crucial?]
Se percorrêssemos $v$ de forma crescente (de $C[i]$ até $k$), o valor atualizado $dp[v]$ em um passo menor poderia ser reaproveitado no mesmo laço para atualizar $dp[v + C[i]]$, permitindo que a moeda $C[i]$ fosse usada infinitas vezes! Ao iterar \textbf{de trás para frente} (de $k$ descendo até $C[i]$), ao calcular $dp[v]$, o valor $dp[v - C[i]]$ ainda pertence estritamente à iteração anterior ($i-1$), garantindo que a moeda $C[i]$ seja contabilizada no máximo uma única vez.
\end{alerta}

\vspace{0.4cm}

% =========================================================================
% QUESTÃO 7
% =========================================================================
\section*{Questão 7: Contagem de Maneiras Distintas de Formar o Troco}

\textbf{Enunciado Oficial:} Dados um conjunto de valores de moedas $C = c_1, c_2, \dots, c_n$ e uma quantia $k$. Escreva um algoritmo que retorne o \textbf{número de maneiras diferentes} que podemos formar a quantia $k$. \emph{Dica: Modifique o algoritmo \textsc{COINS-PD}$(C, k)$.}

\begin{teoria}[Contagem de Partições: Combinações vs. Permutações]
Em problemas de contagem combinatória:
\begin{itemize}
    \item Se a ordem das moedas \textbf{não importa} (combinações / conjuntos, ex: $\{1, 5\}$ é a mesma coisa que $\{5, 1\}$): devemos colocar o laço de moedas \textbf{por fora} e o laço de quantia \textbf{por dentro}.
    \item Se a ordem importasse (permutações, ex: passos de uma escada): o laço de quantia ficaria por fora.
\end{itemize}
Para moedas de dinheiro, o padrão universal em PAA é contar \textbf{combinações distintas}.
Definindo $W[i, v]$ como o número de formas de atingir $v$ usando as moedas $\{c_1, \dots, c_i\}$:
\[
W[i, v] = W[i-1, v] + W[i, v - c_i].
\]
O primeiro termo conta as soluções que não utilizam a moeda $c_i$; o segundo termo conta as soluções que utilizam a moeda $c_i$ ao menos uma vez.
\end{teoria}

\subsection*{Algoritmo em Programação Dinâmica}

\begin{lstlisting}[caption={Algoritmo para Contar Maneiras Distintas de Formar o Troco}]
COUNT-COIN-CHANGE-WAYS(C, n, k)
1  let dp[0 .. k] be a new array
2  dp[0] = 1                         // Ha 1 maneira de formar 0: conjunto vazio
3  for v = 1 to k
4      dp[v] = 0
5  for j = 1 to n                   // Moedas por fora garante combinacoes
6      for v = C[j] to k            // Quantia crescente permite repeticao
7          dp[v] = dp[v] + dp[v - C[j]]
8  return dp[k]
\end{lstlisting}

\begin{dica}[Exemplo Didático Curto]
Seja $C = \{1, 2\}$ e $k = 3$:
\begin{itemize}
    \item Inicial: $dp = [1, 0, 0, 0]$.
    \item Após moeda 1: $dp = [1, 1, 1, 1]$ (apenas $1+1+1$).
    \item Após moeda 2:
    $v=2: dp[2] = dp[2] + dp[0] = 1 + 1 = 2$ ($\{1, 1\}$ e $\{2\}$).
    $v=3: dp[3] = dp[3] + dp[1] = 1 + 1 = 2$ ($\{1, 1, 1\}$ e $\{1, 2\}$).
\end{itemize}
Retorna 2 maneiras. Complexidade: $\Theta(n \cdot k)$ tempo e $\Theta(k)$ espaço.
\end{dica}

\vspace{0.4cm}

% =========================================================================
% QUESTÃO 8
% =========================================================================
\section*{Questão 8: Mochila 0-1 (Knapsack Problem)}

\textbf{Enunciado Oficial:} Mostre a execução passo-a-passo do algoritmo \textsc{KNAPSACK-PD}$(n, K)$ para a seguinte entrada (capacidade $K = 10$):
\begin{center}
\begin{tabular}{c|cc}
\toprule
\textbf{item} $i$ & \textbf{peso} $w_i$ & \textbf{valor} $v_i$ \\
\midrule
1 & 6 & \$30 \\
2 & 3 & \$14 \\
3 & 4 & \$16 \\
4 & 2 & \$9 \\
\bottomrule
\end{tabular}
\end{center}

\begin{teoria}[Subestrutura Ótima e Recorrência da Mochila 0-1]
Para $i \in \{1, \dots, n\}$ e capacidade corrente $k \in \{0, \dots, K\}$:
\[
M[i, k] = 
\begin{cases}
M[i-1, k], & \text{se } w_i > k \\
\max\Big( M[i-1, k], \; M[i-1, k - w_i] + v_i \Big), & \text{se } w_i \le k
\end{cases}
\]
Cada célula $M[i, k]$ representa o valor máximo que cabe em uma mochila de peso $k$ usando um subconjunto dos itens $\{1, \dots, i\}$.
\end{teoria}

\subsection*{Preenchimento Detalhado da Matriz}

\paragraph{Linha $i = 1$ ($w_1 = 6, v_1 = 30$):}
\begin{itemize}
    \item Para $k = 0, 1, 2, 3, 4, 5$: o item não cabe ($w_1 = 6 > k$) $\implies M[1, k] = 0$.
    \item Para $k = 6, 7, 8, 9, 10$: o item cabe $\implies M[1, k] = v_1 = 30$.
\end{itemize}

\paragraph{Linha $i = 2$ ($w_2 = 3, v_2 = 14$):}
\begin{itemize}
    \item $k = 0, 1, 2$: $k < 3 \implies M[2, k] = M[1, k] = 0$.
    \item $k = 3, 4, 5$: $\max(M[1, k], M[1, k-3] + 14) = \max(0, 0 + 14) = 14$.
    \item $k = 6, 7, 8$: $\max(M[1, k], M[1, k-3] + 14) = \max(30, 0 + 14) = 30$.
    \item $k = 9, 10$: $\max(M[1, k], M[1, k-3] + 14) = \max(30, 30 + 14) = 44$.
\end{itemize}

\paragraph{Linha $i = 3$ ($w_3 = 4, v_3 = 16$):}
\begin{itemize}
    \item $k = 0, 1, 2$: $M[3, k] = 0$.
    \item $k = 3$: $k < 4 \implies M[3, 3] = M[2, 3] = 14$.
    \item $k = 4, 5$: $\max(M[2, k], M[2, k-4] + 16) = \max(14, 0 + 16) = 16$.
    \item $k = 6, 7, 8$: $\max(M[2, k], M[2, k-4] + 16) = \max(30, M[2, k-4] + 16) = 30$.
    \item $k = 9$: $\max(M[2, 9], M[2, 5] + 16) = \max(44, 14 + 16 = 30) = 44$.
    \item $k = 10$: $\max(M[2, 10], M[2, 6] + 16) = \max(44, 30 + 16 = \mathbf{46}) = \mathbf{46}$.
\end{itemize}

\paragraph{Linha $i = 4$ ($w_4 = 2, v_4 = 9$):}
\begin{itemize}
    \item $k = 0, 1$: $0$.
    \item $k = 2$: $\max(M[3, 2], M[3, 0] + 9) = \max(0, 9) = 9$.
    \item $k = 3$: $\max(M[3, 3], M[3, 1] + 9) = \max(14, 9) = 14$.
    \item $k = 4$: $\max(M[3, 4], M[3, 2] + 9) = \max(16, 9) = 16$.
    \item $k = 5$: $\max(M[3, 5], M[3, 3] + 9) = \max(16, 14 + 9 = 23) = 23$.
    \item $k = 6, 7$: $\max(30, 16 + 9 = 25) = 30$.
    \item $k = 8$: $\max(30, 30 + 9 = 39) = 39$.
    \item $k = 9$: $\max(44, 30 + 9 = 39) = 44$.
    \item $k = 10$: $\max(M[3, 10], M[3, 8] + 9) = \max(46, 30 + 9 = 39) = \mathbf{46}$.
\end{itemize}

\textbf{Tabela Completa $M[i, k]$:}
\begin{center}
\begin{tabular}{c|ccccccccccc}
\toprule
$i \setminus k$ & 0 & 1 & 2 & 3 & 4 & 5 & 6 & 7 & 8 & 9 & 10 \\
\midrule
1 ($w_1=6, v_1=30$) & 0 & 0 & 0 & 0 & 0 & 0 & \textbf{30} & 30 & 30 & 30 & 30 \\
2 ($w_2=3, v_2=14$) & 0 & 0 & 0 & 14 & 14 & 14 & 30 & 30 & 30 & 44 & 44 \\
3 ($w_3=4, v_3=16$) & 0 & 0 & 0 & 14 & 16 & 16 & 30 & 30 & 30 & 44 & \textbf{46} \\
4 ($w_4=2, v_4=9$)  & 0 & 0 & 9 & 14 & 16 & 23 & 30 & 30 & 39 & 44 & \textbf{46} \\
\bottomrule
\end{tabular}
\end{center}

\subsection*{Reconstrução da Solução Ótima (Backtracking)}
\begin{enumerate}
    \item Célula $M[4, 10] = 46$. Como $M[4, 10] = M[3, 10] = 46$, o valor ótimo veio da linha de cima $\implies$ \textbf{Item 4 NÃO foi escolhido}. Capacidade continua $10$.
    \item Célula $M[3, 10] = 46$. Como $M[3, 10] \ne M[2, 10]$ ($46 \ne 44$), o valor ótimo veio da inclusão do item 3: $M[2, 10 - 4] + 16 = 30 + 16 = 46 \implies$ \textbf{Item 3 FOI ESCOLHIDO}. Nova capacidade: $10 - w_3 = 10 - 4 = 6$.
    \item Célula $M[2, 6] = 30$. Como $M[2, 6] = M[1, 6] = 30$, veio da linha de cima $\implies$ \textbf{Item 2 NÃO foi escolhido}. Capacidade continua $6$.
    \item Célula $M[1, 6] = 30$. Como $M[1, 6] = v_1 = 30 > 0 \implies$ \textbf{Item 1 FOI ESCOLHIDO}. Nova capacidade: $6 - 6 = 0$.
\end{enumerate}

\textbf{Solução Ótima:} Itens $\{1, 3\}$, com peso total $6 + 4 = 10\text{ kg}$ e valor máximo de \textbf{\$46}.

\begin{teoria}[Conceito Crucial: Por que a Mochila é Pseudo-Polinomial?]
A complexidade de tempo da Mochila 0-1 é $\Theta(n \cdot W)$.
\textbf{Pegadinha clássica de prova:} Esse tempo não é polinomial no tamanho da entrada! Em teoria da computação, o tamanho da entrada de um inteiro $W$ é o seu número de bits: $m = \lceil \log_2 W \rceil + 1$.
Portanto:
\[
\Theta(n \cdot W) = \Theta(n \cdot 2^m).
\]
O algoritmo é, na verdade, \textbf{exponencial no comprimento da representação binária de $W$}. Chamamos essa classe de complexidade de \textbf{pseudo-polinomial}.
\end{teoria}

\vspace{0.4cm}

% =========================================================================
% QUESTÃO 9
% =========================================================================
\section*{Questão 9: Mochila com Repetições (Unbounded Knapsack)}

\textbf{Enunciado Oficial:} Escreva um algoritmo que resolve o problema da mochila no qual \textbf{é permitido repetições de itens na mochila}.

\begin{teoria}[Mochila 0-1 vs. Mochila Não Limitada (Unbounded Knapsack)]
A diferença fundamental entre a Mochila 0-1 e a Mochila com Repetições reside na dependência de estados:
\begin{itemize}
    \item \textbf{Mochila 0-1 (Sem repetição):} Ao decidir incluir o item $i$, a capacidade restante $w - w_i$ deve ser resolvida usando apenas os itens anteriores:
    \[
    M[i, w] = \max\Big( M[i-1, w], \; M[\mathbf{i-1}, w - w_i] + v_i \Big).
    \]
    \item \textbf{Mochila Não Limitada (Com repetição):} Ao incluir o item $i$, ele pode ser escolhido novamente na capacidade restante. Logo, o estado recursivo permanece na mesma linha $i$:
    \[
    M[i, w] = \max\Big( M[i-1, w], \; M[\mathbf{i}, w - w_i] + v_i \Big).
    \]
\end{itemize}
Com isso, podemos colapsar toda a matriz em um único vetor unidimensional $K[0 \dots W]$.
\end{teoria}

\subsection*{Algoritmo em Programação Dinâmica}

\begin{lstlisting}[caption={Algoritmo da Mochila com Repetições (Unbounded Knapsack)}]
UNBOUNDED-KNAPSACK(w, v, n, W)
1  let K[0 .. W] be a new array
2  K[0] = 0
3  for c = 1 to W
4      K[c] = 0
5      for i = 1 to n
6          if w[i] <= c
7              K[c] = max(K[c], K[c - w[i]] + v[i])
8  return K[W]
\end{lstlisting}

\begin{dica}[Algoritmo com Reconstrução dos Itens]
Para recuperar quais itens foram colocados na mochila com repetição:
\begin{lstlisting}
UNBOUNDED-KNAPSACK-WITH-ITEMS(w, v, n, W)
1  let K[0 .. W] and choice[0 .. W] be new arrays
2  K[0] = 0
3  for c = 1 to W
4      K[c] = 0
5      choice[c] = -1
6      for i = 1 to n
7          if w[i] <= c and K[c - w[i]] + v[i] > K[c]
8              K[c] = K[c - w[i]] + v[i]
9              choice[c] = i
10 let items_list be an empty list
11 curr = W
12 while curr > 0 and choice[curr] != -1
13     it = choice[curr]
14     add it to items_list
15     curr = curr - w[it]
16 return K[W] and items_list
\end{lstlisting}
\begin{itemize}
    \item \textbf{Complexidade de Tempo:} $\Theta(n \cdot W)$.
    \item \textbf{Complexidade de Espaço:} $O(W)$ (redução expressiva de memória em relação à matriz $n \times W$).
\end{itemize}
\end{dica}

\vspace{0.4cm}

% =========================================================================
% QUESTÃO 10
% =========================================================================
\section*{Questão 10: Subsequência Comum Mais Longa (LCS Completa)}

\textbf{Enunciado Oficial:} Escreva um algoritmo que retorna a LCS de duas sequências (não apenas o tamanho da LCS). \emph{Dica: Modifique o algoritmo \textsc{LCS-LENGTH}$(S, n)$ para registrar as soluções parciais do problema.}

\begin{teoria}[Princípio de Indução da LCS]
Dadas duas sequências $X = \langle x_1, \dots, x_n \rangle$ e $Y = \langle y_1, \dots, y_m \rangle$:
\begin{enumerate}
    \item \textbf{Caso 1 ($x_n = y_m$):} O último caractere coincide necessariamente. Logo, ele faz parte da LCS, e o problema se reduz a encontrar a LCS dos prefixos $X[1 \dots n-1]$ e $Y[1 \dots m-1]$:
    \[
    M[i, j] = M[i-1, j-1] + 1 \quad (\text{movimento diagonal } \nwarrow).
    \]
    \item \textbf{Caso 2 ($x_n \ne y_m$):} Os caracteres finais não coincidem. Pelo menos um deles não pertence à LCS ótima. Portanto:
    \[
    M[i, j] = \max\Big( M[i-1, j], \; M[i, j-1] \Big) \quad (\text{movimento para cima } \uparrow \text{ ou esquerda } \leftarrow).
    \]
\end{enumerate}
\end{teoria}

\subsection*{Abordagem 1: Com Matriz de Setas $B$ (Cormen - CLRS)}

\begin{lstlisting}[caption={Preenchimento com Matriz de Rota $B$}]
LCS-LENGTH-WITH-TABLE(X, n, Y, m)
1  let M[0 .. n, 0 .. m] and B[1 .. n, 1 .. m] be new matrices
2  for i = 0 to n
3      M[i, 0] = 0
4  for j = 1 to m
5      M[0, j] = 0
6  for i = 1 to n
7      for j = 1 to m
8          if X[i] == Y[j]
9              M[i, j] = M[i - 1, j - 1] + 1
10             B[i, j] = "DIAG"
11         else if M[i - 1, j] >= M[i, j - 1]
12             M[i, j] = M[i - 1, j]
13             B[i, j] = "UP"
14         else
15             M[i, j] = M[i, j - 1]
16             B[i, j] = "LEFT"
17 return M and B
\end{lstlisting}

\begin{lstlisting}[caption={Impressão Recursiva da Cadeia de Caracteres}]
PRINT-LCS(B, X, i, j)
1  if i == 0 or j == 0
2      return
3  if B[i, j] == "DIAG"
4      PRINT-LCS(B, X, i - 1, j - 1)
5      print X[i]
6  else if B[i, j] == "UP"
7      PRINT-LCS(B, X, i - 1, j)
8  else
9      PRINT-LCS(B, X, i, j - 1)
\end{lstlisting}

\subsection*{Abordagem 2: Backtracking Direto sem Matriz Auxiliar $B$}
Podemos poupar $\Theta(n \cdot m)$ de memória reconstruindo o caminho diretamente a partir dos valores numéricos de $M$:

\begin{lstlisting}[caption={Reconstrução Iterativa Direta (Economia de Memória)}]
RETRIEVE-LCS-DIRECT(M, X, Y, n, m)
1  let seq be an empty list
2  i = n
3  j = m
4  while i > 0 and j > 0
5      if X[i] == Y[j]
6          add X[i] to seq
7          i = i - 1
8          j = j - 1
9      else if M[i - 1, j] >= M[i, j - 1]
10         i = i - 1
11     else
12         j = j - 1
13 reverse(seq)
14 return seq
\end{lstlisting}

\begin{dica}[Resumo de Complexidades da LCS]
\begin{itemize}
    \item \textbf{Tempo:} $\Theta(n \cdot m)$ no preenchimento da matriz e $O(n + m)$ no retrocesso (pois a cada passo decrementamos $i$, $j$ ou ambos). Total: $\Theta(n \cdot m)$.
    \item \textbf{Espaço:} $\Theta(n \cdot m)$ para a matriz de valores $M$.
\end{itemize}
\end{dica}

\end{document}
